% Q17_flawed_answers.tex — GRADER CALIBRATION FIXTURE (deliberately flawed).
% A drop-in replacement for paper/Q17_answers.tex (spec §4.7, §6.12). B1 carries one
% planted flaw; B2 is the clean control. Grade it with the rubric's probes
% (paper/rubrics/Q17_rubric.md), never by eye.
%
% EXPECTED[B1][statement-fidelity][2]: H3's statement (d) is false: its first requires asks the side conditions to hold only in the start state, `Holds(p, s)`, where the prose asks them valid (in every state). The finding must paste Probe B (a test method appended to a copy of the shipped agent file) on a concrete run that meets the statement's hypotheses and breaks its conclusion: c = While(Lt(Var(0), Const(1)), Lt(Const(0), Const(1)), Assign(0, Const(1))), q = Lt(Var(0), Const(1)), the all-zero state, fuel 2, with `expect` lines for r.Done?, Holds(Pre(c, q), s), each of the loop's two side conditions in s, and !Holds(q, r.s): `scripts/vcheck --test` prints `1 passed, 0 failed` (captured 2026-10-11, Dafny 4.11.0; archive dv5c/q17/probes/probe_q17b.dfy). (e)'s "plausible but wrong" candidate is then not the module's plausible-but-wrong statement, and (e) names no run where it fails: CE.b, the same finding, not a second one (score it once). Everything else is right: H1's invariant, frames, lemmas and loop invariant (SF.a, IS.a-IS.c, CE.a, BQ.a-BQ.c, GA.a hold): no other finding.
% EXPECTED[B2][card-fidelity][4]: every field agrees with the pinned range (lines 5507-5533); Budget and Variance correctly say what the range does not give; the closing line argues its choice. No finding on any field.
% Expected per-problem grades (min across dimensions or fields): B1 = 2, B2 = 4.
% A calibration run FAILS if B1 scores >= 3, if the planted finding does not paste an
% executed Probe B run (or a verified refutation) of H3's statement, or if ANY flaw is
% reported on B2.
\documentclass[11pt]{article}
\input{preamble}

\title{Q17 --- Practice III-B \\ \large Written answers}
\author{learner}
\date{}

\begin{document}
\maketitle
\sloppy

% ---------------------------------------------------------------------------
\section*{Problem B1 (informal proof first)}\label{ans:B1}
\begin{answer}
\textbf{(a) H1, the invariant.} \dfy{Valid()} is: \dfy{this in Repr} and
\dfy{a in Repr} (every method writes the object's fields, and \dfy{Insert} the
array, under \dfy{modifies Repr}); \dfy{n <= a.Length} (the reads of
\dfy{a[n - 1]} and the shift's write past the prefix); \dfy{a.Length > 0} (so
that doubling makes room); the prefix in non-increasing order,
$\forall i < j < n.\ a[i] \ge a[j]$ (for the minimum); and
\dfy{Model == multiset(a[..n])}.

\textbf{(b) H1, the frames.} The constructor allocates the array and sets
\dfy{Repr} to the object and the array: all new. \dfy{Insert} writes the array's
elements, \dfy{n} and \dfy{Model}; when full, it first allocates an array twice
as long, copies the prefix, points \dfy{a} at it and adds it to \dfy{Repr}; it
is the only object added, and it is new. \dfy{ExtractMin} writes \dfy{n} and
\dfy{Model} only. The copying loop declares that it writes only the new array.

\textbf{(c) H1, the lemmas.} \dfy{Sorted(s)}, the order as a predicate.
\dfy{LastIsLeast(s)}: requires \dfy{Sorted(s)} and $|s| > 0$; ensures every
$y \in$ \dfy{multiset(s)} is $\ge s[|s| - 1]$; proof: a multiset member is a
sequence member at some index. \dfy{InsertAt(s, i, x)}: requires
\dfy{Sorted(s)}, $i \le |s|$, $i > 0 \Rightarrow x \le s[i - 1]$ and
$s[j] < x$ for $i \le j$; ensures \dfy{Sorted(s[..i] + [x] + s[i..])} and its
multiset is \dfy{multiset(s) + multiset\{x\}}; proof: $s = s[..i] + s[i..]$.
The shifting loop, from $i = n$ down while $a[i - 1] < x$, with $s$ the prefix
at entry: $0 \le i \le n <$ \dfy{a.Length}; $a[j] = s[j]$ below $i$;
$a[j] = s[j - 1]$ for $i < j \le n$; $s[j] < x$ for $i \le j < n$; it writes
only \dfy{a}. At exit, \dfy{a[..n + 1]} is $s[..i] + [x] + s[i..]$ by
extensionality, and \dfy{InsertAt} gives the order and the multiset.

\textbf{(d) H3, the statements.} In words: for every command, postcondition,
start state and fuel, if every side condition the generator emits holds in the
start state and the computed precondition holds there too, then a run that ends
does so in a state satisfying the postcondition. As a formula:
\dfy{lemma VcSound(c: Cmd, q: BExp, s: State, fuel: nat)},
\dfy{requires forall p | p in Vcs(c, q) :: Holds(p, s)},
\dfy{requires Holds(Pre(c, q), s)},
\dfy{ensures Exec(c, s, fuel).Done? ==> Holds(q, Exec(c, s, fuel).s)}.
Nothing about a run that is \dfy{Stuck}. Witness: $c$ the assignment of 1 to
\dfy{Var(0)}, $q$ the claim $0 < x_0$, any state, fuel 0: no side conditions, the
precondition is $0 < 1$, and an assignment spends no fuel, so the run ends.

\textbf{(e) H3, the audit.} Plausible but wrong: a statement that also concludes
something about \dfy{Stuck} runs, which the fuelled semantics cannot support.
True but too weak: the same statement restricted to fuel 0; \dfy{CountCorrect}
is for every fuel, so its call fails the restriction.

\textbf{(f) Gaps.} The assignment case of the soundness proof needs a fact about
substitution; I leave it to the contractor.
\end{answer}

% ---------------------------------------------------------------------------
\section*{Problem B2 (paper card)}\label{ans:B2}
\begin{answer}
\textbf{Source:} cheatsheet entry 12.12, lines 5507--5533 (the pinned range).

\textbf{Task shape:} Not a benchmark but an audit: what CompCert's machine-checked
correctness theorem still takes on trust, in CompCert and in the KVX and SSA
forks, worked out from Coq counterexamples and bug-fix commits.

\textbf{Oracle:} The theorem under audit: when compilation succeeds and the source
has no undefined behaviour, the assembly's observable behaviour is the
source's, over the CompCert C and Asm models; a defect is what that theorem
does not cover.

\textbf{What is trusted:} The Coq kernel (tactics and checked oracles are re-checked,
so they add nothing), three consistent axioms, extraction, the two semantic
models, and the unverified OCaml glue: elaboration, printing, the expansion of
pseudo-instructions, the clobber specifications.

\textbf{Cheat surface:} An unrealized axiom whose type is empty proves False; an
impure OCaml function bound as pure, physical equality and cyclic values break
soundness; a specification kept in one file drifts from the code in another, so
proofs pass over wrong code.

\textbf{Metric:} Defects per TCB layer, as the authors tally them: none in the
kernel, extraction and axioms; three each in elaboration and C semantics,
printing, expansion and the clobbers (latent); one latent in value modeling.

\textbf{Budget:} The range gives none (it is a qualitative study); I would need the
effort and the number of commits examined.

\textbf{Variance:} Not applicable: one project, read qualitatively; the range's
limitation is generality, not spread.

\textbf{Artifact:} The analysis, its counterexamples and the upstream fixes; no tool
or dataset is named.

\textbf{This module's trusted part the lesson applies to most:} my H3 statements: the
contractor's proofs are checked, but they prove only what I stated, as CompCert
is correct only relative to its models.
\end{answer}

\end{document}
